% Luc Destombe --- licence CC BY-NC-SA 4.0
% https://creativecommons.org/licenses/by-nc-sa/4.0/deed.fr
% Fichier autonome : compiler avec pdflatex, deux fois (signets, nombre de pages).
\def\lyceeversion{1.0}
\documentclass[11pt,a4paper]{article}
\makeatletter
\newif\ifeleve
\elevefalse

\newif\iflycee@palatino
\lycee@palatinotrue

\RequirePackage[utf8]{inputenc}
\RequirePackage[T1]{fontenc}
\RequirePackage[french]{babel}
\RequirePackage{etoolbox}

\newcommand{\ShorthandsOff}{\@ifundefined{shorthandoff}{}{\shorthandoff{;:!?}}}
\newcommand{\ShorthandsOn}{\@ifundefined{shorthandon}{}{\shorthandon{;:!?}}}
\ShorthandsOff
\AtBeginDocument{\ShorthandsOn}
\AtBeginEnvironment{tikzpicture}{\ShorthandsOff}

\RequirePackage{geometry}
\RequirePackage{fancyhdr}
\RequirePackage{lastpage}
\RequirePackage{multicol}
\RequirePackage{array}
\RequirePackage{enumitem}
\RequirePackage{xcolor}
\RequirePackage{needspace}

\RequirePackage{amsmath,amssymb}
\RequirePackage{mathpazo}
\renewcommand{\ttdefault}{lmtt}
\DeclareMathSizes{10.95}{10.95}{8}{5.5}
\begingroup\catcode`\_=\active \catcode`\^=\active
\gdef_#1{\sb{\scriptscriptstyle #1}}
\gdef^#1{\sp{\scriptscriptstyle\lycee@moinsepais #1}}
\endgroup
\newcommand\lycee@moins{\mathbin{%
  \setbox\z@\hbox{$\m@th\scriptscriptstyle\lycee@moinsorig$}%
  \dimen@=.55\wd\z@ \advance\dimen@-.5pt
  \hbox to.75\wd\z@{\hss$\m@th\scriptscriptstyle\vcenter{\hbox{%
    \tikz\draw[line width=.5pt,line cap=round](0,0)--(\the\dimen@,0);}}$\hss}}}
\begingroup\lccode`\~=`\- \lowercase{\endgroup\let~}\lycee@moins
\newcommand\lycee@moinsepais{\mathcode`\-="8000 }
\AtBeginDocument{\mathchardef\lycee@moinsorig=\mathcode`\-\relax}
\AtBeginDocument{\catcode`\_=12 \mathcode`\_="8000
  \catcode`\^=12 \mathcode`\^="8000 }
\newcommand\lycee@exposants{%
  \fontdimen13\textfont2=.48\fontdimen6\textfont2
  \fontdimen14\textfont2=.48\fontdimen6\textfont2
  \fontdimen15\textfont2=.48\fontdimen6\textfont2
  \fontdimen13\scriptfont2=.48\fontdimen6\scriptfont2
  \fontdimen14\scriptfont2=.48\fontdimen6\scriptfont2
  \fontdimen15\scriptfont2=.48\fontdimen6\scriptfont2}
\every@math@size\expandafter{\the\every@math@size\lycee@exposants}
\newlength{\retraitcalculs}
\setlength{\retraitcalculs}{2em}

\RequirePackage{tikz}
\usetikzlibrary{arrows.meta,calc,babel,shapes.misc}
\RequirePackage[most]{tcolorbox}
\tcbuselibrary{skins,breakable}

\definecolor{coulDef}{HTML}{1F4E79}
\definecolor{coulProp}{HTML}{7030A0}
\definecolor{coulRem}{HTML}{C55A11}
\definecolor{coulEx}{HTML}{2E7D32}
\definecolor{coulExo}{HTML}{B03A2E}
\definecolor{coulPart}{HTML}{1F4E79}
\definecolor{coulMeth}{HTML}{1B6E4F}
\definecolor{coulCourbe}{HTML}{0B5ED7}
\definecolor{coulCourbeB}{HTML}{C00000}
\definecolor{coulCourbeC}{HTML}{2E7D32}
\definecolor{coulGrille}{HTML}{8C8C8C}
\definecolor{coulRep}{HTML}{1B6E4F}
\definecolor{coulBar}{HTML}{2E7D32}

\newcommand{\R}{\mathbb{R}}
\newcommand{\Rs}{\mathbb{R}^{*}}
\newcommand{\N}{\mathbb{N}}
\newcommand{\Ns}{\mathbb{N}^{*}}
\newcommand{\Z}{\mathbb{Z}}
\newcommand{\Q}{\mathbb{Q}}
\renewcommand{\le}{\leqslant}
\renewcommand{\ge}{\geqslant}
\renewcommand{\approx}{\simeq}
\newcommand{\vect}[1]{\overrightarrow{#1}}
\newcommand{\norme}[1]{\left\lVert #1 \right\rVert}
\newcommand{\intervalle}[2]{\left[#1\,;#2\right]}
\RequirePackage{cancel}
\renewcommand{\CancelColor}{\color{coulExo}}
\newcommand{\fois}[1]{\times{\color{coulExo}#1}}
\newcommand{\barre}[1]{\cancel{#1}}
\newlength{\lycee@eqh}\newlength{\lycee@eqd}
\newsavebox{\lycee@eqa}\newsavebox{\lycee@eqb}\newsavebox{\lycee@eqc}
\newcommand{\lycee@eqcol}[1]{\hbox{$\renewcommand{\arraystretch}{1.35}%
  \begin{array}[t]{@{}>{\displaystyle}r@{}>{\displaystyle{}}l@{}}#1\end{array}$}}
\newcommand{\lycee@eqmesure}[1]{%
  \ifdim\ht#1>\lycee@eqh\setlength{\lycee@eqh}{\ht#1}\fi
  \ifdim\dp#1>\lycee@eqd\setlength{\lycee@eqd}{\dp#1}\fi}
\newcommand{\lycee@eqtrait}{\hspace{1.2em}\tikz[baseline=0pt]\draw[dash pattern=on 2pt off 2pt,
  line width=0.5pt] (0,\lycee@eqh) -- (0,-\lycee@eqd);\hspace{1.2em}}
\newcommand{\eqproduit}[2]{\par\smallskip\noindent C'est une équation produit :
  $A\times B=0 \iff A=0$ ou $B=0$.\par\nopagebreak\smallskip
  \begingroup\setlength{\lycee@eqh}{0pt}\setlength{\lycee@eqd}{0pt}%
  \setbox\lycee@eqa=\lycee@eqcol{#1}\setbox\lycee@eqb=\lycee@eqcol{#2}%
  \lycee@eqmesure\lycee@eqa\lycee@eqmesure\lycee@eqb%
  \noindent\hspace*{\retraitcalculs}\usebox\lycee@eqa\lycee@eqtrait\usebox\lycee@eqb\endgroup\par\smallskip}
\newcommand{\eqproduittrois}[3]{\par\smallskip\noindent C'est une équation produit :
  $A\times B\times C=0 \iff A=0$ ou $B=0$ ou $C=0$.\par\nopagebreak\smallskip
  \begingroup\setlength{\lycee@eqh}{0pt}\setlength{\lycee@eqd}{0pt}%
  \setbox\lycee@eqa=\lycee@eqcol{#1}\setbox\lycee@eqb=\lycee@eqcol{#2}\setbox\lycee@eqc=\lycee@eqcol{#3}%
  \lycee@eqmesure\lycee@eqa\lycee@eqmesure\lycee@eqb\lycee@eqmesure\lycee@eqc%
  \noindent\hspace*{\retraitcalculs}\usebox\lycee@eqa\lycee@eqtrait\usebox\lycee@eqb\lycee@eqtrait\usebox\lycee@eqc\endgroup
  \par\smallskip}

\tikzset{grille/.style={coulGrille,line width=0.4pt}}
\newcommand{\quadrillage}[4]{\draw[grille,step=1] (#1,#2) grid (#3,#4);}
\tikzset{
  croix/.style={cross out,draw,line width=0.6pt,inner sep=0pt,minimum size=4pt},
  croixdroite/.style={croix,rotate=45,line width=0.8pt},
}
\newcommand{\pt}[3]{\path (#1) node[croix]{} node[#2,font=\small,inner sep=2.5pt]{$#3$};}

\newcommand{\lycee@repere}[4]{%
  \draw[grille,step=1] (#1,#2) grid (#3,#4);
  \draw[-{Stealth[length=2mm]},thick] (#1,0)--({#3+0.4},0) node[below left=-1pt]{$x$};
  \draw[-{Stealth[length=2mm]},thick] (0,#2)--(0,{#4+0.4}) node[below left=-1pt]{$y$};
  \node[below left,font=\scriptsize,inner sep=1.5pt] at (0,0) {$0$};}
\newcommand{\lycee@gradx}[1]{\foreach \k/\l in {#1}{%
  \draw (\k,0.07)--(\k,-0.07) node[below,font=\scriptsize,inner sep=1.5pt]{$\l$};}}
\newcommand{\lycee@grady}[1]{\foreach \k/\l in {#1}{%
  \draw (0.07,\k)--(-0.07,\k) node[left,font=\scriptsize,inner sep=1.5pt]{$\l$};}}
\newcommand{\lycee@intff}[2]{\left[#1\,;#2\right]}
\newcommand{\lycee@intfo}[2]{\left[#1\,;#2\right[}
\newcommand{\lycee@intof}[2]{\left]#1\,;#2\right]}
\newcommand{\lycee@intoo}[2]{\left]#1\,;#2\right[}
\newcommand{\lycee@ens}[1]{\left\{#1\right\}}
\AtBeginDocument{%
  \providecommand{\repere}{\lycee@repere}%
  \providecommand{\gradx}{\lycee@gradx}%
  \providecommand{\grady}{\lycee@grady}%
  \providecommand{\Cf}{\mathcal{C}\sb{\scriptscriptstyle f}}%
  \providecommand{\Cg}{\mathcal{C}\sb{\scriptscriptstyle g}}%
  \providecommand{\intff}{\lycee@intff}%
  \providecommand{\intfo}{\lycee@intfo}%
  \providecommand{\intof}{\lycee@intof}%
  \providecommand{\intoo}{\lycee@intoo}%
  \providecommand{\ens}{\lycee@ens}}

\newlist{questions}{enumerate}{3}
\setlist[questions,1]{label=\arabic*\textdegree),leftmargin=*,itemsep=2pt,topsep=3pt}
\setlist[questions,2]{label=\alph*),leftmargin=*,itemsep=1pt,topsep=2pt}
\setlist[questions,3]{label=\roman*),leftmargin=*,itemsep=1pt,topsep=1pt}
\setlist[itemize]{label=\textbullet,leftmargin=*,itemsep=2pt,topsep=3pt}

\newcommand{\besoinplace}[1]{\par\penalty\@M\begingroup
  \dimen@=#1\relax\dimen@ii=\pagegoal\advance\dimen@ii-\pagetotal
  \ifdim\dimen@>\dimen@ii\ifdim\dimen@ii>\z@\vfil\fi\break\fi\endgroup}

\newcommand{\lycee@pied}{\thepage/\pageref{LastPage}}
\newcommand{\entete}[2]{%
  \pagestyle{fancy}\fancyhf{}%
  \fancyhead[L]{\footnotesize\color{coulPart}\textbf{#1}}%
  \fancyhead[R]{\footnotesize\color{coulPart}\textbf{#2}}%
  \fancyfoot[C]{\footnotesize\lycee@pied}%
  \renewcommand{\headrulewidth}{0.4pt}%
  \renewcommand{\headrule}{\hbox to\headwidth{%
    \color{coulPart!40}\leaders\hrule height \headrulewidth\hfill}}}

\RequirePackage{ccicons}
\newcommand{\lycee@licence}{{\scriptsize\color{gray}%
  \copyright\ Luc Destombe --- \ccbyncsa\ CC BY-NC-SA 4.0}}
\newif\iflycee@licence \lycee@licencetrue
\newcommand{\sanslicence}{\lycee@licencefalse}
\AtBeginDocument{\iflycee@licence\AddToHookNext{shipout/foreground}{%
  \put(0,0){\rlap{\kern\dimexpr1in+\hoffset+\oddsidemargin+\textwidth\relax
    \llap{\raisebox{-\dimexpr1in+\voffset+\topmargin+\headheight+\headsep
      +\textheight+\footskip\relax}{\lycee@licence}}}}}\fi}

\newcommand{\lycee@site}{faire-des-maths.fr}
\newcommand{\lycee@version}{{\scriptsize\color{gray}\lycee@site
  \ifdefined\lyceeversion\ --- Version~\lyceeversion\fi}}
\AtBeginDocument{\iflycee@licence\AddToHookNext{shipout/foreground}{%
  \put(0,0){\rlap{\kern\dimexpr1in+\hoffset+\oddsidemargin\relax
    \raisebox{-\dimexpr1in+\voffset+\topmargin+\headheight+\headsep
      +\textheight+\footskip\relax}{\lycee@version}}}}\fi}

\newcommand{\formulesserrees}{%
  \setlength{\abovedisplayskip}{3pt}\setlength{\belowdisplayskip}{3pt}%
  \setlength{\abovedisplayshortskip}{2pt}\setlength{\belowdisplayshortskip}{3pt}}

\setlength{\parindent}{0pt}

\RequirePackage[implicit=false,hidelinks,pdfencoding=auto,psdextra,bookmarksopen,
  bookmarksopenlevel=1,pdfcreator={},pdfproducer={}]{hyperref}
\RequirePackage{bookmark}
\hypersetup{pdfauthor={Luc Destombe},pdfkeywords={CC BY-NC-SA 4.0}}
\newcounter{lycee@signet}
\newcommand{\signet}[3][]{\stepcounter{lycee@signet}%
  \edef\lycee@dest{\if\relax\detokenize{#1}\relax
    signet.\arabic{lycee@signet}\else#1\fi}%
  \hypertarget{\lycee@dest}{}\bookmark[level=#2,dest=\lycee@dest]{#3}}
\pdfstringdefDisableCommands{%
  \def\R{R}\def\N{N}\def\Z{Z}\def\Q{Q}\def\vect#1{#1}\def\Cf{Cf}\def\Cg{Cg}\def\fois#1{×#1}%
  \def\barre#1{#1}\def\intervalle#1#2{[#1;#2]}\def\intff#1#2{[#1;#2]}%
  \def\textsuperscript#1{#1}\def\dfrac#1#2{#1/#2}\def\frac#1#2{#1/#2}%
  \def\vec#1{#1⃗}}%
\newcommand{\res}[1]{#1}
\geometry{top=2cm,bottom=1cm,left=1cm,right=1cm,headheight=14pt,footskip=0.6cm}

\RequirePackage{pgfplots}
\pgfplotsset{compat=1.18}
\RequirePackage{tkz-tab}

\newcounter{partiecours}
\renewcommand{\thepartiecours}{\Roman{partiecours}}
\newcommand{\partiecours}[1]{%
  \refstepcounter{partiecours}%
  \Needspace*{8\baselineskip}%
  \bigskip
  \noindent\begin{tcolorbox}[enhanced,colback=coulPart,colframe=coulPart,
      boxrule=0pt,arc=2pt,left=8pt,right=8pt,top=4pt,bottom=4pt,
      before skip=6pt,after skip=10pt]
    \color{white}\bfseries\large \thepartiecours{} --- #1%
    \signet[partie-\arabic{partiecours}]{1}{\thepartiecours{} --- #1}
  \end{tcolorbox}%
  \addcontentsline{toc}{section}{\thepartiecours{} --- #1}%
}

\newtcolorbox{boitecours}[3][]{%
  enhanced,breakable,
  colback=white,colframe=#2,coltitle=white,
  fonttitle=\bfseries,
  attach boxed title to top left={xshift=8pt,yshift=-\tcboxedtitleheight/2},
  boxed title style={colback=#2,arc=2pt,boxrule=0pt},
  boxrule=0.9pt,arc=2pt,
  left=8pt,right=8pt,top=8pt,bottom=6pt,
  title={#3},#1}

\newcounter{methode}
\newenvironment{definitionbox}[1][Définition]%
  {\begin{boitecours}[phantom={\signet{2}{#1}}]{coulDef}{#1}}{\end{boitecours}}
\newenvironment{proprietebox}[1][Propriété]%
  {\begin{boitecours}[phantom={\signet{2}{#1}}]{coulProp}{#1}}{\end{boitecours}}
\newenvironment{methodebox}[1][Méthode]{\stepcounter{methode}%
  \begin{boitecours}[phantom={\signet[methode-\arabic{methode}]{2}{#1}}]{coulMeth}{#1}}%
  {\end{boitecours}}

\newcommand{\etape}[1]{\par\smallskip\textbf{\color{coulMeth}#1}\par\nobreak\smallskip}
\newcommand{\enonce}[1]{\emph{#1}\par\smallskip}
\newcommand{\reponse}[1]{\par\smallskip
  {\footnotesize\itshape\color{black!55}Réponses : #1}}
\newcommand{\demo}[1]{\textbf{\color{coulMeth}Démonstration.} #1}
\newcommand{\afaire}[1]{%
  \par\smallskip
  \noindent{\small\color{coulExo}$\blacktriangleright$\ \textbf{À faire :}
    fiche d'exercices du chapitre, #1.}\par\smallskip}

\newenvironment{calculs}{\setlength{\jot}{5pt}\@fleqntrue\@mathmargin\retraitcalculs
  \setlength{\abovedisplayskip}{4pt}\setlength{\belowdisplayskip}{4pt}%
  \setlength{\abovedisplayshortskip}{4pt}\setlength{\belowdisplayshortskip}{4pt}%
  \csname alignat*\endcsname{2}}%
  {\csname endalignat*\endcsname}
\newcommand{\rem}[1]{\text{\small\itshape\color{black!60}#1}}
\newcommand{\explique}[2]{\par\smallskip\noindent
  \begin{minipage}[t]{0.57\linewidth}\raggedright #1\end{minipage}\hfill
  \begin{minipage}[t]{0.40\linewidth}\raggedright\leavevmode
    {\small\itshape\color{black!60}#2}\end{minipage}\par}
\newcommand{\cas}[1]{\par\smallskip\noindent\textbf{\color{coulExo}#1}\vadjust{\nobreak}}
\newcommand{\calc}[1]{\par\smallskip\textbf{\color{coulExo}#1}\par\nobreak}

\newtcolorbox{remarquebox}[1][Remarque]{%
  enhanced,breakable,colback=white,colframe=coulRem,
  boxrule=0pt,leftrule=2.5pt,arc=0pt,outer arc=0pt,
  left=8pt,right=6pt,top=5pt,bottom=5pt,
  before upper={\textbf{\color{coulRem}#1}\par\smallskip}}

\newtcolorbox{exemplebox}[1][Exemple]{%
  enhanced,breakable,colback=white,colframe=coulEx,
  boxrule=0pt,leftrule=2.5pt,arc=0pt,outer arc=0pt,
  left=8pt,right=6pt,top=5pt,bottom=5pt,
  before upper={\textbf{\color{coulEx}#1}\par\smallskip}}

\pgfplotsset{
  repere/.style={
    axis lines=middle,
    axis line style={-{Stealth[length=2.4mm]},thick},
    every axis x label/.style={at={(ticklabel* cs:1.0)},anchor=west},
    every axis y label/.style={at={(ticklabel* cs:1.0)},anchor=south},
    label style={font=\footnotesize},
    tick label style={font=\scriptsize},
    grid=both,
    major grid style={grille},
    minor tick num=0,
    tick align=inside,
    clip mode=individual,
  },
}
\tikzset{
  courbe/.style={coulCourbe,line width=1.1pt,smooth},
  courbeB/.style={coulCourbeB,line width=1.1pt,smooth},
  courbeC/.style={coulCourbeC,line width=1.1pt,smooth},
}

\newlist{etapes}{enumerate}{1}
\setlist[etapes]{label=\textbf{\arabic*.},leftmargin=*,itemsep=1pt,topsep=2pt}

\newlist{expressions}{enumerate}{1}
\setlist[expressions]{label={\Alph*\ $=$},leftmargin=*,itemsep=3pt,topsep=3pt}

\newcounter{exo}
\renewcommand{\theexo}{\ifnum\value{exo}<10 0\fi\arabic{exo}}
\newtcolorbox{exercice}[1][]{%
  enhanced,breakable,
  colback=white,colframe=coulExo,
  boxrule=0.9pt,arc=2pt,
  left=8pt,right=8pt,top=8pt,bottom=6pt,
  coltitle=white,fonttitle=\bfseries,
  attach boxed title to top left={xshift=8pt,yshift=-\tcboxedtitleheight/2},
  boxed title style={colback=coulExo,arc=2pt,boxrule=0pt},
  phantom={\signet{2}{Exercice \arabic{exo}}},
  title={Exercice \theexo\ifx\relax#1\relax\else{} --- #1\fi}}
\newenvironment{exo}[1][\relax]{\refstepcounter{exo}\begin{exercice}[#1]}{\end{exercice}}

  \newenvironment{correction}{\par}{\par}
  \newcommand{\autableau}{}
  \newcommand{\etapeexemples}{\etape{Exemples corrigés}}
  \newcommand{\etapetableau}[1]{\etape{#1}}
  \newenvironment{exempletableau}[1][Exemple]%
    {\begin{exemplebox}[#1]}{\end{exemplebox}}

\RequirePackage{comment}
\excludecomment{correctionavj}

\newcommand{\coursEntete}{}
\newcommand{\coursNiveau}{}
\newcommand{\coursTitre}{}
\newcommand{\coursSurTitre}{}
\newcommand{\coursSousTitre}{}

\newcommand{\enteteCours}[1]{\renewcommand{\coursEntete}{#1}}
\newcommand{\chapitrecours}[2]{\enteteCours{Chapitre #1 --- #2}}
\newcommand{\niveaucours}[1]{\renewcommand{\coursNiveau}{#1}}
\newcommand{\titrecours}[3][]{%
  \renewcommand{\coursTitre}{#2}%
  \renewcommand{\coursSurTitre}{#1}%
  \renewcommand{\coursSousTitre}{#3}}

\pagestyle{fancy}
\fancyhf{}
\fancyhead[L]{\footnotesize\color{coulPart}\textbf{\coursEntete}}
\fancyhead[R]{\footnotesize\color{coulPart}\textbf{\coursNiveau}}
\fancyfoot[C]{\footnotesize\thepage}
\renewcommand{\headrulewidth}{0.4pt}
\renewcommand{\headrule}{\hbox to\headwidth{\color{coulPart!40}\leaders\hrule height \headrulewidth\hfill}}

\setlength{\parskip}{2pt}

\newcommand{\lycee@mentioneleve}{\\[2pt]{\small\itshape Version élève}}
\AtBeginDocument{%
  \begin{center}
    \begin{tcolorbox}[enhanced,width=0.72\linewidth,colback=white,colframe=coulPart,
        boxrule=1.6pt,arc=3pt,halign=center,top=10pt,bottom=10pt]
      \ifx\coursSurTitre\empty
        {\Huge\bfseries\color{coulPart} \coursTitre}\\[4pt]
      \else
        {\Huge\bfseries\color{coulPart} \coursTitre}\\[3pt]
        {\Large\bfseries\color{coulPart} \coursSurTitre}\\[5pt]
      \fi
      {\normalsize \coursSousTitre}%
      
    \end{tcolorbox}
  \end{center}%
}

\RequirePackage{tikz}
\RequirePackage{xcolor}

\usetikzlibrary{decorations.text}
\definecolor{lycee@metal}{HTML}{1F4E79}
\definecolor{lycee@manche}{HTML}{8A5A2B}
\definecolor{lycee@etincelle}{HTML}{F28C28}
\newcommand{\lycee@pluriel}[2]{#1 #2\ifnum#1>1 s\fi}
\newcommand{\lycee@sceau}{%
  \edef\lycee@bas{\lycee@pluriel{\lyceerelectures}{relecture} et
    \lycee@pluriel{\lyceerodages}{rodage}}%
  \begin{tikzpicture}[line join=round]
    \fill[coulPart!10,draw=coulPart,line width=.8pt] (0:2.2)
      \foreach \i in {1,...,71}{ -- (\i*5:{isodd(\i)?2.1:2.2}) } -- cycle;
    \draw[coulPart,line width=1.2pt] (0,0) circle (1.93);
    \draw[coulPart,line width=.4pt] (0,0) circle (1.85);
    \draw[coulPart,line width=.4pt] (0,0) circle (1.2);
    \path[decorate,decoration={text along path,text align=center,
        text={|\fontsize{9.5}{11}\selectfont\bfseries\color{coulPart}|TEST{\'E} AVEC MES {\'E}L{\`E}VES}}]
      (180:1.45) arc (180:0:1.45);
    \path[decorate,decoration={text along path,text align=center,
        text/.expanded={|\noexpand\fontsize{8.5}{10}\noexpand\selectfont\noexpand\color{coulPart}|\lycee@bas}}]
      (195:1.68) arc (195:345:1.68);
    \begin{scope}[x=.25mm,y=-.25mm]
      \begin{scope}[shift={(-32,-40)}]%
        \fill[lycee@metal] (23,31) rectangle (41,49);
        \draw[lycee@etincelle,line width=.5pt,line cap=round]
          (20,30)--(15,27) (19,34)--(14,35) (44,30)--(49,27) (45,34)--(50,35);
        \fill[lycee@manche,rounded corners=.3mm] (38,24.6) rectangle (60,27.6);
        \fill[lycee@metal,rounded corners=.4mm] (25,20) rectangle (39,31);
        \fill[lycee@etincelle,rounded corners=.2mm] (25,20) rectangle (28,31);
      \end{scope}
    \end{scope}
    \ifdefined\lyceesceauplusieurs\def\lycee@annee{Ann{\'e}es~:}\else
      \def\lycee@annee{Ann{\'e}e}\fi
    \path[decorate,decoration={text along path,text align=center,
        text/.expanded={|\noexpand\fontsize{8}{10}\noexpand\selectfont\noexpand\color{coulPart}|%
          \unexpanded\expandafter{\lycee@annee} \unexpanded\expandafter{\lyceesceauannees}}}]
      (95:1.08) arc (95:445:1.08);
  \end{tikzpicture}}
\AtEndDocument{\ifdefined\lyceesceauannees
  \par\besoinplace{5cm}\vspace{1.5em}%
  \begin{center}\lycee@sceau\end{center}\fi}
\makeatother
\tcbuselibrary{listings}

\chapitrecours{11}{Probabilités}
\niveaucours{Seconde}
\titrecours{PROBABILITÉS}{Cours --- Seconde}

\begin{document}

\newcommand{\Card}{\operatorname{Card}}

\newtcblisting{python}{%
  listing only,enhanced,
  colback=black!3,colframe=black!35,boxrule=0.5pt,arc=1pt,
  left=4pt,right=4pt,top=1pt,bottom=1pt,
  listing options={language=Python,basicstyle=\ttfamily\small,
    keywordstyle=\color{coulMeth}\bfseries,columns=fullflexible,
    keepspaces=true,showstringspaces=false,aboveskip=0pt,belowskip=0pt}}

\tikzset{
  arbre/.style={grow'=right,level distance=2.3cm,
    level 1/.style={sibling distance=2.2cm},
    level 2/.style={sibling distance=1.1cm},
    every node/.style={font=\small}},
  poids/.style={font=\scriptsize,sloped,inner sep=1.5pt},
}

\partiecours{Vocabulaire des probabilités}

On lance un dé, on tire une carte, on choisit un élève au hasard : on connaît
les résultats possibles, mais on ne peut pas prévoir celui qui va sortir.

\begin{definitionbox}[Définition 1 --- Expérience aléatoire, issue, univers]
\begin{itemize}
  \item Une expérience est \textbf{aléatoire} quand on connaît tous ses résultats
    possibles sans pouvoir prévoir lequel se produira.
  \item Chaque résultat possible est une \textbf{issue}.
  \item L'ensemble de toutes les issues est l'\textbf{univers}, noté $\Omega$
    (« oméga ») ; son nombre d'éléments se note $\Card(\Omega)$ (« cardinal de
    $\Omega$ »).
\end{itemize}
Avec un dé à six faces : $\Omega=\ens{1\,;2\,;3\,;4\,;5\,;6}$ et $\Card(\Omega)=6$.
\end{definitionbox}

\begin{definitionbox}[Définition 2 --- Événement]
\begin{itemize}
  \item Un \textbf{événement} est un ensemble d'issues : on le décrit par une phrase,
    ou par la liste de ses issues entre accolades.
  \item Il est \textbf{réalisé} quand l'issue obtenue en fait partie.
  \item Un événement qui ne contient aucune issue est \textbf{impossible} ; il se note
    $\varnothing$ (« ensemble vide »). Un événement qui contient toutes les issues
    est \textbf{certain} : c'est $\Omega$.
\end{itemize}
Avec un dé : $A$ : « obtenir un nombre pair » s'écrit $A=\ens{2\,;4\,;6}$. Si le dé
donne $4$, l'événement $A$ est réalisé.
\end{definitionbox}

\begin{methodebox}[Méthode 1 --- Décrire l'univers et des événements]
\etapeexemples
Une application de quiz tire au hasard une question parmi $15$ questions numérotées
de $1$ à $15$, et on note son numéro.
\begin{questions}[label=\alph*)]
  \item Écrire l'univers $\Omega$, puis donner $\Card(\Omega)$.
  \item Écrire sous forme d'ensemble les événements $A$ : « le numéro est un multiple
    de $4$ » et $B$ : « le numéro est supérieur ou égal à $12$ ».
  \item Écrire l'événement $C$ : « le numéro est un multiple de $20$ ». Comment
    s'appelle un tel événement ?
  \item La question n°$12$ sort. Lesquels des événements $A$ et $B$ sont réalisés ?
\end{questions}
\begin{correction}
\cas{a)}
\explique{$\Omega=\res{\ens{1\,;2\,;3\,;\dots\,;15}}$ et $\Card(\Omega)=\res{15}$.}{les
$15$ numéros possibles}
\cas{b)}
\explique{$A=\res{\ens{4\,;8\,;12}}$.}{$4\times 4=16$ dépasse $15$}
\explique{$B=\res{\ens{12\,;13\,;14\,;15}}$.}{« supérieur ou égal » : $12$ en fait
partie}
\cas{c)}
\explique{Aucun numéro de $1$ à $15$ n'est un multiple de $20$, donc
$C=\res{\varnothing}$ : c'est un \res{événement impossible}.}{aucune issue}
\cas{d)}
\explique{$12$ appartient à $A$ et à $B$ : \res{les événements $A$ et $B$ sont
réalisés tous les deux}.}{on cherche $12$ dans chaque liste}
\end{correction}

\etape{À vous de jouer}
On lance un dé équilibré à $12$ faces numérotées de $1$ à $12$, et on note le numéro
de la face du dessus.
\begin{questions}[label=\alph*)]
  \item Écrire l'univers $\Omega$, puis donner $\Card(\Omega)$.
  \item Écrire sous forme d'ensemble les événements $A$ : « le numéro est un multiple
    de $5$ » et $B$ : « le numéro est inférieur ou égal à $3$ ».
  \item Écrire l'événement $C$ : « le numéro est supérieur ou égal à $13$ ». Comment
    s'appelle un tel événement ?
  \item Le dé donne $10$. Lesquels des événements $A$ et $B$ sont réalisés ?
\end{questions}
\reponse{a) $\Omega=\ens{1\,;2\,;\dots\,;12}$, $\Card(\Omega)=12$ ; b)~$A=\ens{5\,;10}$,
$B=\ens{1\,;2\,;3}$ ; c)~$C=\varnothing$, événement impossible ; d) seul $A$.}
\end{methodebox}

\afaire{exercices 1 à 3}

\partiecours{Calcul d'une probabilité}

\begin{proprietebox}[Propriété 1 --- Une probabilité est un nombre entre $0$ et $1$]
La probabilité d'un événement $A$ se note $P(A)$ ; elle mesure les chances que $A$ se
réalise.
\begin{itemize}
  \item Pour tout événement $A$ : $0\le P(A)\le 1$.
  \item $P(\varnothing)=0$ (événement impossible) et $P(\Omega)=1$ (événement certain).
\end{itemize}
$P(A)=0{,}8$ signifie que $A$ a $8$ chances sur $10$ de se réaliser, soit $80\,\%$.
\end{proprietebox}

\begin{definitionbox}[Définition 3 --- Équiprobabilité]
Quand toutes les issues ont la même probabilité, on parle d'\textbf{équiprobabilité}.
C'est un choix de modèle, qu'on ne démontre pas : les mots « au hasard », « dé
équilibré », « jetons indiscernables au toucher » indiquent ce choix.
\end{definitionbox}

\begin{proprietebox}[Propriété 2 --- Probabilité en situation d'équiprobabilité]
\[
P(A)=\frac{\Card(A)}{\Card(\Omega)}
    =\frac{\text{nombre d'issues de }A}{\text{nombre total d'issues}}
\]
\end{proprietebox}

\begin{methodebox}[Méthode 2 --- Calculer une probabilité en situation d'équiprobabilité]
\etapeexemples
\begin{questions}[label=\alph*)]
  \item On lance un dé équilibré à $8$ faces numérotées de $1$ à $8$. Calculer les
    probabilités des événements $A$ : « obtenir un nombre pair », $B$ : « obtenir un
    multiple de $3$ » et $C$ : « obtenir un nombre inférieur à $9$ ».
  \item Un sac contient $5$ jetons rouges, $7$ jetons verts et $3$ jetons bleus,
    indiscernables au toucher. On tire un jeton au hasard. Calculer les probabilités
    des événements $V$ : « le jeton est vert » et $N$ : « le jeton n'est pas rouge ».
\end{questions}
\begin{correction}
\cas{a)}
\explique{Le dé est équilibré : les $8$ issues ont la même probabilité, et
$\Card(\Omega)=8$.}{on justifie l'équiprobabilité}
\explique{$A=\ens{2\,;4\,;6\,;8}$, donc $\Card(A)=4$.}{on liste les issues de $A$}
\begin{calculs}
P(A) &= \dfrac{\Card(A)}{\Card(\Omega)} &\qquad& \rem{propriété 2}\\
P(A) &= \dfrac{4}{8} &&\\
P(A) &= \dfrac{1\times\barre{4}}{2\times\barre{4}} && \rem{on simplifie par $4$}\\
P(A) &= \res{\dfrac{1}{2}} &&
\end{calculs}
\explique{$B=\ens{3\,;6}$, donc $P(B)=\dfrac{2}{8}=\res{\dfrac{1}{4}}$.}{même calcul,
plus court}
\explique{Tous les numéros sont inférieurs à $9$, donc $C=\Omega$ et
$P(C)=\res{1}$.}{événement certain}
\cas{b)}
\explique{Les jetons sont indiscernables : les issues sont les $15$ jetons, qui ont
tous la même probabilité.}{piège : les $3$ couleurs, elles, n'ont pas la même
probabilité}
\explique{$7$ jetons sur $15$ sont verts, donc $P(V)=\res{\dfrac{7}{15}}$.}{$\Card(V)=7$}
\explique{Les jetons qui ne sont pas rouges sont les $7$ verts et les $3$ bleus, soit
$10$ jetons.}{$\Card(N)=7+3$}
\begin{calculs}
P(N) &= \dfrac{10}{15} &\qquad& \rem{$10$ issues sur $15$}\\
P(N) &= \dfrac{2\times\barre{5}}{3\times\barre{5}} && \rem{on simplifie par $5$}\\
P(N) &= \res{\dfrac{2}{3}} &&
\end{calculs}
\end{correction}

\etape{À vous de jouer}
\begin{questions}[label=\alph*)]
  \item On lance un dé équilibré à $10$ faces numérotées de $1$ à $10$. Calculer les
    probabilités des événements $A$ : « obtenir un nombre pair », $B$ : « obtenir un
    multiple de $4$ » et $C$ : « obtenir un nombre inférieur ou égal à $10$ ».
  \item Un sac contient $4$ billes jaunes, $9$ billes noires et $2$ billes blanches,
    indiscernables au toucher. On tire une bille au hasard. Calculer les probabilités
    des événements $N$ : « la bille est noire » et $J$ : « la bille n'est pas jaune ».
\end{questions}
\reponse{a) $P(A)=\frac{1}{2}$, $P(B)=\frac{1}{5}$, $P(C)=1$ ; b)~$P(N)=\frac{3}{5}$,
$P(J)=\frac{11}{15}$.}
\end{methodebox}

\afaire{exercices 4 à 7}

\partiecours{Intersection, réunion, événement contraire}

On retrouve l'intersection, la réunion et le complémentaire vus au chapitre 2 avec
les intervalles.

\begin{definitionbox}[Définition 4 --- Intersection, réunion, contraire]
$A$ et $B$ sont deux événements d'un même univers $\Omega$.
\begin{itemize}
  \item $A\cap B$ (« $A$ inter $B$ ») : les issues qui sont dans $A$ \textbf{et}
    dans $B$ ; $A$ et $B$ se réalisent en même temps.
  \item $A\cup B$ (« $A$ union $B$ ») : les issues qui sont dans $A$ \textbf{ou} dans
    $B$, ou dans les deux ; au moins l'un des deux se réalise.
  \item $\overline{A}$ (« $A$ barre ») : les issues de $\Omega$ qui ne sont
    \textbf{pas} dans $A$ ; c'est l'événement \textbf{contraire} de $A$.
  \item $A$ et $B$ sont \textbf{incompatibles} quand $A\cap B=\varnothing$ : ils ne
    peuvent pas se réaliser en même temps.
\end{itemize}
\begin{center}
\begin{tikzpicture}[x=1cm,y=1cm]
  \begin{scope}
    \clip (-0.35,0) circle (0.6);
    \fill[coulMeth!30] (0.35,0) circle (0.6);
  \end{scope}
  \draw (-1.25,-0.8) rectangle (1.25,0.8);
  \draw (-0.35,0) circle (0.6) (0.35,0) circle (0.6);
  \node[font=\scriptsize] at (-0.65,0) {$A$};
  \node[font=\scriptsize] at (0.65,0) {$B$};
  \node[font=\scriptsize] at (1.05,0.6) {$\Omega$};
  \node[font=\small] at (0,-1.1) {$A\cap B$};
\end{tikzpicture}
\qquad
\begin{tikzpicture}[x=1cm,y=1cm]
  \fill[coulMeth!30] (-0.35,0) circle (0.6) (0.35,0) circle (0.6);
  \draw (-1.25,-0.8) rectangle (1.25,0.8);
  \draw (-0.35,0) circle (0.6) (0.35,0) circle (0.6);
  \node[font=\scriptsize] at (-0.65,0) {$A$};
  \node[font=\scriptsize] at (0.65,0) {$B$};
  \node[font=\scriptsize] at (1.05,0.6) {$\Omega$};
  \node[font=\small] at (0,-1.1) {$A\cup B$};
\end{tikzpicture}
\qquad
\begin{tikzpicture}[x=1cm,y=1cm]
  \fill[coulMeth!30] (-1.25,-0.8) rectangle (1.25,0.8);
  \fill[white] (-0.35,0) circle (0.6);
  \draw (-1.25,-0.8) rectangle (1.25,0.8);
  \draw (-0.35,0) circle (0.6) (0.35,0) circle (0.6);
  \node[font=\scriptsize] at (-0.65,0) {$A$};
  \node[font=\scriptsize] at (0.65,0) {$B$};
  \node[font=\scriptsize] at (1.05,0.6) {$\Omega$};
  \node[font=\small] at (0,-1.1) {$\overline{A}$};
\end{tikzpicture}
\end{center}
\end{definitionbox}

\begin{methodebox}[Méthode 3 --- Décrire $A\cap B$, $A\cup B$ et $\overline{A}$]
\etapeexemples
\begin{questions}[label=\alph*)]
  \item Une application de quiz tire au hasard une question parmi $12$ questions
    numérotées de $1$ à $12$. On note $A$ : « le numéro est un multiple de $3$ » et
    $B$ : « le numéro est supérieur ou égal à $9$ ». Écrire $A$, $B$, $A\cap B$,
    $A\cup B$ et $\overline{A}$ sous forme d'ensembles. $A$ et $B$ sont-ils
    incompatibles ?
  \item On choisit au hasard un élève d'un lycée. On note $S$ : « l'élève est en
    seconde » et $D$ : « l'élève est demi-pensionnaire ». Décrire par une phrase les
    événements $S\cap D$, $S\cup D$ et $\overline{S}\cap D$.
\end{questions}
\begin{correction}
\cas{a)}
\explique{$A=\ens{3\,;6\,;9\,;12}$ et $B=\ens{9\,;10\,;11\,;12}$.}{on liste d'abord
$A$ et $B$}
\explique{$A\cap B=\res{\ens{9\,;12}}$.}{les numéros communs aux deux listes}
\explique{$A\cup B=\res{\ens{3\,;6\,;9\,;10\,;11\,;12}}$.}{les numéros d'au moins une
liste, écrits une seule fois}
\explique{$\overline{A}=\res{\ens{1\,;2\,;4\,;5\,;7\,;8\,;10\,;11}}$.}{les numéros qui ne
sont pas dans $A$}
\explique{$A\cap B\ne\varnothing$, donc \res{$A$ et $B$ ne sont pas
incompatibles}.}{$9$ et $12$ sont dans les deux}
\cas{b)}
\explique{$S\cap D$ : \res{« l'élève est un demi-pensionnaire de seconde »}.}{« et » :
les deux à la fois}
\explique{$S\cup D$ : \res{« l'élève est en seconde ou demi-pensionnaire »}, ou les
deux.}{« ou » : au moins l'un des deux}
\explique{$\overline{S}\cap D$ : \res{« l'élève est demi-pensionnaire et n'est pas
en seconde »}.}{$\overline{S}$ : « l'élève n'est pas en seconde »}
\end{correction}

\etape{À vous de jouer}
\begin{questions}[label=\alph*)]
  \item On lance un dé équilibré à $12$ faces numérotées de $1$ à $12$. On note
    $A$ : « le numéro est impair » et $B$ : « le numéro est inférieur ou égal à
    $4$ ». Écrire $A$, $B$, $A\cap B$, $A\cup B$ et $\overline{A}$ sous forme
    d'ensembles. $A$ et $B$ sont-ils incompatibles ?
  \item On choisit au hasard un spectateur d'un festival. On note $M$ : « le
    spectateur a moins de $18$ ans » et $T$ : « le spectateur est venu en train ».
    Décrire par une phrase les événements $M\cap T$, $M\cup T$ et $\overline{M}\cap T$.
\end{questions}
\reponse{a) $A\cap B=\ens{1\,;3}$, $A\cup B=\ens{1\,;2\,;3\,;4\,;5\,;7\,;9\,;11}$,
$\overline{A}=\ens{2\,;4\,;6\,;8\,;10\,;12}$ ; pas incompatibles ; b) voir la correction.}
\end{methodebox}

\afaire{exercices 8 et 9}

\begin{proprietebox}[Propriété 3 --- Probabilité du contraire et d'une réunion]
Pour tous événements $A$ et $B$ :
\begin{center}
$P\big(\overline{A}\big)=1-P(A)$ \qquad $P(A\cup B)=P(A)+P(B)-P(A\cap B)$
\end{center}
Dans $P(A)+P(B)$, les issues de $A\cap B$ sont comptées deux fois : on retire
$P(A\cap B)$ une fois. Si $A$ et $B$ sont incompatibles, $P(A\cup B)=P(A)+P(B)$.
\end{proprietebox}

\begin{methodebox}[Méthode 4 --- Calculer $P\big(\overline{A}\big)$ et $P(A\cup B)$]
\etapeexemples
\begin{questions}[label=\alph*)]
  \item Une application de quiz tire au hasard une question parmi $12$ questions
    numérotées de $1$ à $12$. On note $A$ : « le numéro est un multiple de $3$ » et
    $B$ : « le numéro est supérieur ou égal à $9$ ». Calculer $P(A)$, $P(B)$ et
    $P(A\cap B)$, puis $P\big(\overline{A}\big)$ et $P(A\cup B)$ avec la propriété 3.
  \item Sur un site de vente en ligne, la probabilité qu'un client achète un livre
    est $0{,}3$, celle qu'il achète un jeu est $0{,}45$, et celle qu'il achète les deux
    est $0{,}1$. On note $L$ et $J$ ces deux événements. Calculer
    $P\big(\overline{L}\big)$, $P(L\cup J)$, puis la probabilité qu'un client n'achète
    ni livre ni jeu.
\end{questions}
\begin{correction}
\cas{a)}
\explique{$A=\ens{3\,;6\,;9\,;12}$, $B=\ens{9\,;10\,;11\,;12}$ et
$A\cap B=\ens{9\,;12}$.}{méthode 3}
\explique{Équiprobabilité : $P(A)=\dfrac{4}{12}=\res{\dfrac{1}{3}}$ et
$P(B)=\dfrac{4}{12}=\res{\dfrac{1}{3}}$.}{$\Card(\Omega)=12$}
\explique{$P(A\cap B)=\dfrac{2}{12}=\res{\dfrac{1}{6}}$.}{$2$ issues}
\begin{calculs}
P\big(\overline{A}\big) &= 1-P(A) &\qquad& \rem{probabilité du contraire}\\
P\big(\overline{A}\big) &= \dfrac{3}{3}-\dfrac{1}{3} && \rem{$1=\frac{3}{3}$}\\
P\big(\overline{A}\big) &= \res{\dfrac{2}{3}} &&
\end{calculs}
\begin{calculs}
P(A\cup B) &= P(A)+P(B)-P(A\cap B) &\qquad& \rem{propriété 3}\\
P(A\cup B) &= \dfrac{1}{3}+\dfrac{1}{3}-\dfrac{1}{6} &&\\
P(A\cup B) &= \dfrac{1\fois{2}}{3\fois{2}}+\dfrac{1\fois{2}}{3\fois{2}}-\dfrac{1}{6}
  && \rem{même dénominateur $6$}\\
P(A\cup B) &= \dfrac{2+2-1}{6} &&\\
P(A\cup B) &= \dfrac{3}{6} &&\\
P(A\cup B) &= \res{\dfrac{1}{2}} &&
\end{calculs}
\explique{Vérification : $A\cup B=\ens{3\,;6\,;9\,;10\,;11\,;12}$ compte $6$ issues, et
$\dfrac{6}{12}=\dfrac{1}{2}$.}{on compte les issues}
\cas{b)}
\begin{calculs}
P\big(\overline{L}\big) &= 1-P(L) &\qquad& \rem{probabilité du contraire}\\
P\big(\overline{L}\big) &= 1-0{,}3 &&\\
P\big(\overline{L}\big) &= \res{0{,}7} &&
\end{calculs}
\begin{calculs}
P(L\cup J) &= P(L)+P(J)-P(L\cap J) &\qquad& \rem{propriété 3}\\
P(L\cup J) &= 0{,}3+0{,}45-0{,}1 &&\\
P(L\cup J) &= \res{0{,}65} &&
\end{calculs}
\explique{« Ni livre ni jeu » est le contraire de $L\cup J$ (au moins l'un des deux).}{on
traduit la phrase}
\begin{calculs}
P\big(\overline{L\cup J}\big) &= 1-P(L\cup J) &\qquad& \rem{probabilité du contraire}\\
P\big(\overline{L\cup J}\big) &= 1-0{,}65 &&\\
P\big(\overline{L\cup J}\big) &= \res{0{,}35} &&
\end{calculs}
\end{correction}

\etape{À vous de jouer}
\begin{questions}[label=\alph*)]
  \item On lance un dé équilibré à $12$ faces numérotées de $1$ à $12$. On note
    $A$ : « le numéro est impair » et $B$ : « le numéro est inférieur ou égal à
    $4$ ». Calculer $P(A)$, $P(B)$ et $P(A\cap B)$, puis $P\big(\overline{B}\big)$ et
    $P(A\cup B)$ avec la propriété 3.
  \item Dans une médiathèque, la probabilité qu'un inscrit emprunte une bande dessinée
    est $0{,}4$, celle qu'il emprunte un film est $0{,}25$, et celle qu'il emprunte les
    deux est $0{,}15$. On note $D$ et $F$ ces deux événements. Calculer
    $P\big(\overline{D}\big)$, $P(D\cup F)$, puis la probabilité qu'un inscrit
    n'emprunte ni bande dessinée ni film.
\end{questions}
\reponse{a) $P(A)=\frac{1}{2}$, $P(B)=\frac{1}{3}$, $P(A\cap B)=\frac{1}{6}$,
$P\big(\overline{B}\big)=\frac{2}{3}$, $P(A\cup B)=\frac{2}{3}$ ; b)~$0{,}6$ ; $0{,}5$ ;
$0{,}5$.}
\end{methodebox}

\afaire{exercices 8, 10 et 11}

\partiecours{Loi de probabilité}

\begin{definitionbox}[Définition 5 --- Loi de probabilité]
Donner la \textbf{loi de probabilité} d'une expérience, c'est donner la probabilité de
chaque issue, souvent dans un tableau. Ces probabilités sont positives ou nulles, et
leur somme vaut $1$.

\smallskip
Une urne contient $2$ boules rouges, $3$ vertes et $5$ bleues ; on tire une boule au
hasard et on note sa couleur :
\begin{center}
\renewcommand{\arraystretch}{1.5}
\begin{tabular}{|l|c|c|c|}
\hline
Couleur & rouge & verte & bleue \\
\hline
Probabilité & $0{,}2$ & $0{,}3$ & $0{,}5$ \\
\hline
\end{tabular}
\end{center}
Ici, les issues (les couleurs) n'ont pas la même probabilité : on sort du modèle
d'équiprobabilité, et la formule $\dfrac{\Card(A)}{\Card(\Omega)}$ ne s'applique plus.
C'est le cas le plus souvent quand une expérience est donnée par sa loi.
\end{definitionbox}

\begin{proprietebox}[Propriété 4 --- Probabilité d'un événement]
La probabilité d'un événement est la somme des probabilités des issues qui le
composent.
\end{proprietebox}

\begin{remarquebox}[D'où vient la loi ?]
Soit d'un choix d'équiprobabilité (dé équilibré, tirage au hasard), soit de fréquences
observées sur un très grand nombre d'essais (dé truqué, punaise qui tombe sur la tête
ou sur le côté) : c'est la loi des grands nombres de la partie VIII.
\end{remarquebox}

\begin{methodebox}[Méthode 5 --- Compléter une loi de probabilité et l'utiliser]
\etapeexemples
\begin{questions}[label=\alph*)]
  \item On lance un dé truqué à six faces, dont la loi de probabilité est la suivante.
    \begin{center}
    \renewcommand{\arraystretch}{1.5}
    \begin{tabular}{|l|c|c|c|c|c|c|}
    \hline
    Face & $1$ & $2$ & $3$ & $4$ & $5$ & $6$ \\
    \hline
    Probabilité & $0{,}1$ & $0{,}15$ & $\dots$ & $0{,}2$ & $0{,}1$ & $0{,}25$ \\
    \hline
    \end{tabular}
    \end{center}
    Calculer la probabilité d'obtenir $3$, puis les probabilités des événements $A$ :
    « obtenir un nombre pair » et $\overline{A}$.
  \item À la tombola d'une fête de quartier, chaque ticket donne « rien », « une
    boisson » ou « une peluche », avec la loi suivante.
    \begin{center}
    \renewcommand{\arraystretch}{2.4}
    \begin{tabular}{|l|c|c|c|}
    \hline
    Lot & rien & une boisson & une peluche \\
    \hline
    Probabilité & $\dfrac{1}{2}$ & $\dfrac{1}{3}$ & $\dots$ \\
    \hline
    \end{tabular}
    \end{center}
    Calculer la probabilité de gagner une peluche, puis celle de l'événement $G$ :
    « le ticket donne un lot ».
\end{questions}
\begin{correction}
\cas{a)}
\explique{La somme des probabilités vaut $1$.}{définition 5}
\begin{calculs}
P(3) &= 1-(0{,}1+0{,}15+0{,}2+0{,}1+0{,}25) &\qquad& \rem{les cinq autres faces}\\
P(3) &= 1-0{,}8 &&\\
P(3) &= \res{0{,}2} &&
\end{calculs}
\explique{$A=\ens{2\,;4\,;6}$.}{les faces paires}
\begin{calculs}
P(A) &= P(2)+P(4)+P(6) &\qquad& \rem{propriété 4}\\
P(A) &= 0{,}15+0{,}2+0{,}25 &&\\
P(A) &= \res{0{,}6} &&
\end{calculs}
\explique{$P\big(\overline{A}\big)=1-0{,}6=\res{0{,}4}$.}{probabilité du contraire}
\cas{b)}
\begin{calculs}
P(\text{peluche}) &= 1-\dfrac{1}{2}-\dfrac{1}{3} &\qquad& \rem{somme des probabilités : $1$}\\
P(\text{peluche}) &= \dfrac{1\fois{6}}{1\fois{6}}-\dfrac{1\fois{3}}{2\fois{3}}
  -\dfrac{1\fois{2}}{3\fois{2}} && \rem{même dénominateur $6$}\\
P(\text{peluche}) &= \dfrac{6-3-2}{6} &&\\
P(\text{peluche}) &= \res{\dfrac{1}{6}} &&
\end{calculs}
\explique{$G$ est réalisé avec une boisson ou une peluche.}{les issues de $G$}
\begin{calculs}
P(G) &= P(\text{boisson})+P(\text{peluche}) &\qquad& \rem{propriété 4}\\
P(G) &= \dfrac{1\fois{2}}{3\fois{2}}+\dfrac{1}{6} && \rem{même dénominateur $6$}\\
P(G) &= \dfrac{3}{6} &&\\
P(G) &= \res{\dfrac{1}{2}} &&
\end{calculs}
\explique{Vérification : $G=\overline{\text{rien}}$, et
$1-\dfrac{1}{2}=\dfrac{1}{2}$.}{contraire de « rien »}
\end{correction}

\etape{À vous de jouer}
\begin{questions}[label=\alph*)]
  \item On lance un dé truqué à six faces, dont la loi de probabilité est la suivante.
    \begin{center}
    \renewcommand{\arraystretch}{1.5}
    \begin{tabular}{|l|c|c|c|c|c|c|}
    \hline
    Face & $1$ & $2$ & $3$ & $4$ & $5$ & $6$ \\
    \hline
    Probabilité & $0{,}25$ & $0{,}1$ & $0{,}15$ & $\dots$ & $0{,}2$ & $0{,}05$ \\
    \hline
    \end{tabular}
    \end{center}
    Calculer la probabilité d'obtenir $4$, puis les probabilités des événements $A$ :
    « obtenir un nombre pair » et $\overline{A}$.
  \item À la tombola d'un club, chaque ticket donne « rien », « un stylo » ou « une
    casquette », avec la loi suivante.
    \begin{center}
    \renewcommand{\arraystretch}{2.4}
    \begin{tabular}{|l|c|c|c|}
    \hline
    Lot & rien & un stylo & une casquette \\
    \hline
    Probabilité & $\dfrac{2}{3}$ & $\dfrac{1}{4}$ & $\dots$ \\
    \hline
    \end{tabular}
    \end{center}
    Calculer la probabilité de gagner une casquette, puis celle de l'événement $G$ :
    « le ticket donne un lot ».
\end{questions}
\reponse{a) $P(4)=0{,}25$ ; $P(A)=0{,}4$ ; $P\big(\overline{A}\big)=0{,}6$ ; b)
$P(\text{casquette})=\frac{1}{12}$ ; $P(G)=\frac{1}{3}$.}
\end{methodebox}

\afaire{exercices 12 à 15}

\partiecours{Tableaux croisés d'effectifs}

\begin{definitionbox}[Définition 6 --- Tableau croisé d'effectifs]
Un \textbf{tableau croisé} répartit une population selon deux caractères (une ligne
par valeur du premier, une colonne par valeur du second), avec une ligne et une
colonne de totaux. Chaque case donne l'effectif des individus qui ont les deux valeurs
à la fois.
\end{definitionbox}

\begin{proprietebox}[Propriété 5 --- Probabilités lues dans un tableau]
Quand on choisit \textbf{au hasard} un individu de la population, il y a
équiprobabilité :
\[
P(A)=\frac{\text{effectif de }A}{\text{effectif total}}
\]
L'effectif de $A\cap B$ se lit dans une case du tableau, celui de $A$ dans un total.
\end{proprietebox}

\begin{remarquebox}[Rappel du chapitre 7]
Prendre $25\,\%$ d'un effectif, c'est le multiplier par $0{,}25$ ; ainsi $25\,\%$ de $180$
vaut $0{,}25\times 180$, soit $45$. On remplit de même les cases d'un tableau
à partir de pourcentages.
\end{remarquebox}

\begin{methodebox}[Méthode 6 --- Compléter un tableau croisé et calculer des probabilités]
\etapeexemples
Un lycée compte $300$ élèves de seconde ; $180$ sont demi-pensionnaires, les autres
sont externes. $25\,\%$ des demi-pensionnaires et $30\,\%$ des externes font partie de
l'association sportive (AS).
\begin{questions}[label=\alph*)]
  \item Compléter le tableau croisé des effectifs ci-dessous.
    \begin{center}
    \renewcommand{\arraystretch}{1.5}
    \begin{tabular}{|l|c|c|c|}
    \hline
     & À l'AS & Pas à l'AS & Total \\
    \hline
    Demi-pensionnaires & & & \\
    \hline
    Externes & & & \\
    \hline
    Total & & & $300$ \\
    \hline
    \end{tabular}
    \end{center}
  \item On choisit au hasard un élève de seconde. On note $D$ : « l'élève est
    demi-pensionnaire » et $S$ : « l'élève fait partie de l'AS ». Calculer $P(D)$,
    $P(S)$ et $P(D\cap S)$.
  \item Calculer $P(D\cup S)$, puis la probabilité que l'élève soit externe et ne
    fasse pas partie de l'AS.
\end{questions}
\begin{correction}
\cas{a)}
\begin{calculs}
\text{externes} &= 300-180 &\qquad& \rem{les autres élèves}\\
\text{externes} &= 120 &&
\end{calculs}
\begin{calculs}
\text{demi-pensionnaires à l'AS} &= 0{,}25\times 180 &\qquad& \rem{$25\,\%$ de $180$}\\
\text{demi-pensionnaires à l'AS} &= 45 &&
\end{calculs}
\begin{calculs}
\text{externes à l'AS} &= 0{,}3\times 120 &\qquad& \rem{$30\,\%$ de $120$}\\
\text{externes à l'AS} &= 36 &&
\end{calculs}
\explique{Les autres cases s'obtiennent par soustraction ($180-45=135$,
$120-36=84$), puis les totaux des colonnes par addition.}{chaque ligne a pour somme son
total}
\begin{center}
\renewcommand{\arraystretch}{1.5}
\begin{tabular}{|l|c|c|c|}
\hline
 & À l'AS & Pas à l'AS & Total \\
\hline
Demi-pensionnaires & $\res{45}$ & $135$ & $180$ \\
\hline
Externes & $36$ & $84$ & $120$ \\
\hline
Total & $81$ & $219$ & $300$ \\
\hline
\end{tabular}
\end{center}
\cas{b)}
\explique{L'élève est choisi au hasard : équiprobabilité sur les $300$ élèves.}{propriété 5}
\explique{$P(D)=\dfrac{180}{300}=\res{0{,}6}$.}{total de la ligne}
\explique{$P(S)=\dfrac{81}{300}=\res{0{,}27}$.}{total de la colonne}
\explique{$P(D\cap S)=\dfrac{45}{300}=\res{0{,}15}$.}{la case commune}
\cas{c)}
\begin{calculs}
P(D\cup S) &= P(D)+P(S)-P(D\cap S) &\qquad& \rem{propriété 3}\\
P(D\cup S) &= 0{,}6+0{,}27-0{,}15 &&\\
P(D\cup S) &= \res{0{,}72} &&
\end{calculs}
\explique{« Externe et pas à l'AS » : $P\big(\overline{D}\cap\overline{S}\big)
=\dfrac{84}{300}=\res{0{,}28}$.}{la case « externes, pas à l'AS »}
\explique{Vérification : cet événement est le contraire de $D\cup S$, et
$1-0{,}72=0{,}28$.}{ni l'un ni l'autre}
\end{correction}

\etape{À vous de jouer}
Un club de sport compte $400$ adhérents ; $240$ sont adultes, les autres sont jeunes.
$20\,\%$ des adultes et $45\,\%$ des jeunes font de la compétition.
\begin{questions}[label=\alph*)]
  \item Dresser le tableau croisé des effectifs (lignes : adultes, jeunes ; colonnes :
    compétition, pas de compétition ; totaux).
  \item On choisit au hasard un adhérent. On note $A$ : « l'adhérent est adulte » et
    $C$ : « l'adhérent fait de la compétition ». Calculer $P(A)$, $P(C)$ et
    $P(A\cap C)$.
  \item Calculer $P(A\cup C)$, puis la probabilité que l'adhérent soit jeune et ne
    fasse pas de compétition.
\end{questions}
\reponse{a) adultes : $48$ et $192$ ; jeunes : $72$ et $88$ ; totaux $120$ et $280$ ;
b)~$P(A)=0{,}6$, $P(C)=0{,}3$, $P(A\cap C)=0{,}12$ ; c)~$0{,}78$ ; $0{,}22$.}
\end{methodebox}

\afaire{exercices 16 à 18}

\partiecours{Probabilités conditionnelles}

Parfois, on choisit au hasard non pas dans toute la population, mais dans une partie
seulement : « parmi les demi-pensionnaires », « parmi les personnes testées
positives ».

\begin{exemplebox}
Sur les $300$ élèves de seconde d'un lycée, $180$ sont demi-pensionnaires, et $45$
de ces demi-pensionnaires font partie de l'association sportive (AS).
\begin{itemize}
  \item On choisit au hasard un élève de seconde : la probabilité qu'il soit un
    demi-pensionnaire de l'AS vaut $\dfrac{45}{300}$, soit $0{,}15$.
  \item On choisit au hasard un demi-pensionnaire : il ne reste que $180$ élèves
    possibles, et la probabilité qu'il fasse partie de l'AS vaut $\dfrac{45}{180}$,
    soit $0{,}25$.
\end{itemize}
Savoir que l'élève est demi-pensionnaire change le dénominateur : on divise par
$180$, pas par $300$.
\end{exemplebox}

\begin{definitionbox}[Définition 7 --- Probabilité de $B$ sachant $A$]
$A$ est un événement de probabilité non nulle. La probabilité que $B$ se réalise,
\textbf{sachant que} $A$ est réalisé, se note $P_A(B)$. Pour un individu choisi au
hasard dans une population :
\[
P_A(B)=\frac{\Card(A\cap B)}{\Card(A)}
\]
On ne regarde que la ligne (ou la colonne) de $A$ dans le tableau croisé.
\end{definitionbox}

\begin{proprietebox}[Propriété 6 --- Calcul avec les probabilités]
\[
P_A(B)=\frac{P(A\cap B)}{P(A)}
\]
\enonce{Exemple : sur $300$ élèves, $180$ sont demi-pensionnaires, dont $45$ à l'AS.
On a $\dfrac{45}{180}=0{,}25$ et $\dfrac{0{,}15}{0{,}6}=0{,}25$ : les deux calculs
donnent le même résultat.}
\demo{on divise le numérateur et le dénominateur par $\Card(\Omega)$.}
\begin{calculs}
P_A(B) &= \dfrac{\Card(A\cap B)}{\Card(A)} &\qquad& \rem{définition 7}\\
P_A(B) &= \dfrac{\Card(A\cap B)\div\Card(\Omega)}{\Card(A)\div\Card(\Omega)}
  && \rem{même diviseur en haut et en bas}\\
P_A(B) &= \dfrac{P(A\cap B)}{P(A)} && \rem{propriété 2}
\end{calculs}
Conclusion : $P_A(B)=\dfrac{P(A\cap B)}{P(A)}$.
\end{proprietebox}

\begin{remarquebox}[Attention à l'ordre]
En général, $P_A(B)\ne P_B(A)$ : la part des demi-pensionnaires qui font de l'AS
n'est pas la part des élèves de l'AS qui sont demi-pensionnaires.
\end{remarquebox}

\begin{definitionbox}[Définition 8 --- Vocabulaire d'un test de dépistage]
On note $M$ : « la personne est malade » et $T$ : « le test est positif ».
\begin{itemize}
  \item Un \textbf{faux positif} est un test positif chez une personne qui n'est pas
    malade ; un \textbf{faux négatif}, un test négatif chez une personne malade.
  \item La \textbf{sensibilité} du test est la part des malades qu'il détecte (test
    positif) : c'est $P_M(T)$. Sensibilité $0{,}95$ : sur $100$ malades, $95$ ont un
    test positif.
  \item Sa \textbf{spécificité} est la part des non-malades qu'il reconnaît (test
    négatif) : c'est $P_{\overline{M}}\big(\overline{T}\big)$. Spécificité $0{,}9$ :
    sur $100$ personnes non malades, $90$ ont un test négatif.
\end{itemize}
Les deux se lisent dans la ligne de la population concernée (malades, non-malades),
jamais dans la colonne des tests positifs.
\end{definitionbox}

\begin{methodebox}[Méthode 7 --- Calculer une probabilité conditionnelle avec un tableau]
\etapeexemples
\begin{questions}[label=\alph*)]
  \item Le tableau répartit les $300$ élèves de seconde d'un lycée. On note $D$ :
    « l'élève est demi-pensionnaire » et $S$ : « l'élève fait partie de l'AS ».
    \begin{center}
    \renewcommand{\arraystretch}{1.5}
    \begin{tabular}{|l|c|c|c|}
    \hline
     & À l'AS & Pas à l'AS & Total \\
    \hline
    Demi-pensionnaires & $45$ & $135$ & $180$ \\
    \hline
    Externes & $36$ & $84$ & $120$ \\
    \hline
    Total & $81$ & $219$ & $300$ \\
    \hline
    \end{tabular}
    \end{center}
    On choisit au hasard un demi-pensionnaire : calculer la probabilité qu'il fasse
    partie de l'AS. On choisit au hasard un élève de l'AS : calculer la probabilité
    qu'il soit demi-pensionnaire. Noter ces probabilités.
  \item Un test rapide de dépistage d'une maladie a été essayé sur $10\,000$
    personnes. On note $M$ : « la personne est malade » et $T$ : « le test est
    positif ».
    \begin{center}
    \renewcommand{\arraystretch}{1.5}
    \begin{tabular}{|l|c|c|c|}
    \hline
     & Test positif & Test négatif & Total \\
    \hline
    Malades & $95$ & $5$ & $100$ \\
    \hline
    Non malades & $495$ & $9\,405$ & $9\,900$ \\
    \hline
    Total & $590$ & $9\,410$ & $10\,000$ \\
    \hline
    \end{tabular}
    \end{center}
    Rappel : la \emph{sensibilité} est la part des malades dont le test est positif ;
    la \emph{spécificité}, la part des non-malades dont le test est négatif ; un
    \emph{faux positif}, un test positif chez une personne non malade.

    Calculer la sensibilité et la spécificité du test, puis la probabilité qu'une
    personne testée positive soit malade, et enfin la part des faux positifs parmi
    les tests positifs (au centième près).
\end{questions}
\begin{correction}
\cas{a)}
\explique{« Parmi les demi-pensionnaires » : on ne regarde que leur ligne, $45$ élèves
de l'AS sur $180$.}{la condition est $D$}
\explique{$P_D(S)=\dfrac{45}{180}=\res{0{,}25}$.}{$\dfrac{\Card(D\cap S)}{\Card(D)}$ ;
on retrouve les $25\,\%$}
\explique{« Parmi les élèves de l'AS » : on ne regarde que leur colonne, $45$
demi-pensionnaires sur $81$.}{la condition est $S$}
\explique{$P_S(D)=\dfrac{45}{81}=\res{\dfrac{5}{9}}$, soit environ $0{,}56$.}{on
simplifie par $9$}
\explique{$P_D(S)\ne P_S(D)$ : l'ordre compte.}{remarque « Attention à l'ordre »}
\cas{b)}
\explique{Sensibilité : $P_M(T)=\dfrac{95}{100}=\res{0{,}95}$.}{ligne des malades : le
test détecte $95$ malades sur $100$}
\explique{Spécificité :
$P_{\overline{M}}\big(\overline{T}\big)=\dfrac{9\,405}{9\,900}=\res{0{,}95}$.}{ligne des
non-malades : $95\,\%$ ont un test négatif}
\explique{Personne testée positive : on ne regarde que la colonne des $590$ tests
positifs.}{la condition est $T$}
\explique{$P_T(M)=\dfrac{95}{590}\approx\res{0{,}16}$.}{calculatrice}
\explique{Faux positifs : $P_T\big(\overline{M}\big)=\dfrac{495}{590}\approx\res{0{,}84}$.}{non
malades parmi les positifs}
\explique{Le test détecte $95\,\%$ des malades, mais $84\,\%$ des tests positifs sont
de faux positifs : la maladie est rare, et les non-malades sont si nombreux que
$5\,\%$ d'entre eux pèsent lourd. Un test positif doit être confirmé.}{$P_M(T)$ et
$P_T(M)$ sont très différents}
\end{correction}

\etape{À vous de jouer}
\begin{questions}[label=\alph*)]
  \item Le tableau répartit les $400$ adhérents d'un club de sport. On note $A$ :
    « l'adhérent est adulte » et $C$ : « l'adhérent fait de la compétition ».
    \begin{center}
    \renewcommand{\arraystretch}{1.5}
    \begin{tabular}{|l|c|c|c|}
    \hline
     & Compétition & Pas de compétition & Total \\
    \hline
    Adultes & $48$ & $192$ & $240$ \\
    \hline
    Jeunes & $72$ & $88$ & $160$ \\
    \hline
    Total & $120$ & $280$ & $400$ \\
    \hline
    \end{tabular}
    \end{center}
    On choisit au hasard un adulte : calculer la probabilité qu'il fasse de la
    compétition. On choisit au hasard un compétiteur : calculer la probabilité qu'il
    soit adulte. Noter ces probabilités.
  \item Un filtre anti-spam a trié $2\,000$ courriels. On note $S$ : « le courriel est
    un spam » et $I$ : « le filtre le classe comme indésirable ».
    \begin{center}
    \renewcommand{\arraystretch}{1.5}
    \begin{tabular}{|l|c|c|c|}
    \hline
     & Classé indésirable & Non classé & Total \\
    \hline
    Spams & $380$ & $20$ & $400$ \\
    \hline
    Autres courriels & $80$ & $1\,520$ & $1\,600$ \\
    \hline
    Total & $460$ & $1\,540$ & $2\,000$ \\
    \hline
    \end{tabular}
    \end{center}
    Rappel : la \emph{sensibilité} du filtre est la part des spams qu'il classe
    indésirables ; sa \emph{spécificité}, la part des autres courriels qu'il ne classe
    pas ; un \emph{faux positif}, un courriel utile classé indésirable.

    Calculer la sensibilité et la spécificité du filtre, puis la probabilité qu'un
    courriel classé indésirable soit un spam, et enfin la part des faux positifs parmi
    les courriels classés indésirables (au centième près).
\end{questions}
\reponse{a) $P_A(C)=0{,}2$ ; $P_C(A)=0{,}4$ ; b)~$P_S(I)=0{,}95$ ; $P_{\overline{S}}\big(\overline{I}\big)=0{,}95$ ;
$P_I(S)\approx 0{,}83$ ;
$P_I\big(\overline{S}\big)\approx 0{,}17$.}
\end{methodebox}

\afaire{exercices 19 à 22}

\partiecours{Arbres pondérés}

\begin{minipage}[c]{0.56\linewidth}\raggedright
\begin{definitionbox}[Définition 9 --- Arbre pondéré]
Un \textbf{arbre pondéré} représente une expérience en deux étapes. Sur chaque
branche, on écrit sa probabilité :
\begin{itemize}
  \item au premier niveau, $P(A)$ et $P\big(\overline{A}\big)$ ;
  \item au second niveau, des probabilités conditionnelles : après $A$, on écrit
    $P_A(B)$ et $P_A\big(\overline{B}\big)$.
\end{itemize}
Un \textbf{chemin} va du départ à une extrémité : le chemin $A$ puis $B$ représente
$A\cap B$.
\end{definitionbox}
\end{minipage}\hfill
\begin{minipage}[c]{0.42\linewidth}\centering
\begin{tikzpicture}[arbre,level distance=2.5cm]
\node {}
  child {node {$A$}
    child {node {$B$} edge from parent node[poids,above] {$P_A(B)$}}
    child {node {$\overline{B}$} edge from parent node[poids,below] {$P_A(\overline{B})$}}
    edge from parent node[poids,above] {$P(A)$}}
  child {node {$\overline{A}$}
    child {node {$B$} edge from parent node[poids,above] {$P_{\overline{A}}(B)$}}
    child {node {$\overline{B}$} edge from parent node[poids,below] {$P_{\overline{A}}(\overline{B})$}}
    edge from parent node[poids,below] {$P(\overline{A})$}};
\end{tikzpicture}
\end{minipage}

\begin{proprietebox}[Propriété 7 --- Règles de calcul sur un arbre]
\begin{itemize}
  \item La somme des probabilités des branches issues d'un même nœud vaut $1$.
  \item La probabilité d'un chemin est le produit des probabilités de ses branches :
    $P(A\cap B)=P(A)\times P_A(B)$ (c'est la propriété 6, multipliée par $P(A)$).
  \item La probabilité d'un événement est la somme des probabilités des chemins qui
    le réalisent.
\end{itemize}
\end{proprietebox}

\begin{methodebox}[Méthode 8 --- Construire un arbre pondéré et l'utiliser]
\etapeexemples
\begin{questions}[label=\alph*)]
  \item Dans un lycée, $40\,\%$ des élèves viennent en bus, les autres à pied ou à
    vélo. $15\,\%$ des élèves qui viennent en bus arrivent en retard au moins une fois
    par mois, contre $5\,\%$ des autres. On choisit un élève au hasard ; on note $B$ :
    « l'élève vient en bus » et $R$ : « l'élève arrive en retard au moins une fois
    par mois ». Traduire les données par des probabilités, construire l'arbre
    pondéré, puis calculer $P(B\cap R)$ et $P(R)$.
  \item Un sac contient $4$ jetons rouges et $2$ jetons jaunes, indiscernables au
    toucher. On tire un jeton au hasard, on le garde, puis on tire un second jeton.
    Construire l'arbre pondéré, puis calculer la probabilité d'obtenir deux jetons
    rouges, et celle d'obtenir deux jetons de la même couleur.
\end{questions}
\begin{correction}
\cas{a)}
\explique{$P(B)=0{,}4$, $P_B(R)=0{,}15$ et $P_{\overline{B}}(R)=0{,}05$.}{« parmi ceux
qui viennent en bus » : probabilité sachant $B$}
\explique{On complète chaque nœud à $1$, d'où $P\big(\overline{B}\big)=0{,}6$,
$P_B\big(\overline{R}\big)=0{,}85$, $P_{\overline{B}}\big(\overline{R}\big)=0{,}95$.}{somme
des branches d'un nœud : $1$}
\begin{center}
\begin{tikzpicture}[arbre]
\node {}
  child {node {$B$}
    child {node {$R$} edge from parent node[poids,above] {$0{,}15$}}
    child {node {$\overline{R}$} edge from parent node[poids,below] {$0{,}85$}}
    edge from parent node[poids,above] {$0{,}4$}}
  child {node {$\overline{B}$}
    child {node {$R$} edge from parent node[poids,above] {$0{,}05$}}
    child {node {$\overline{R}$} edge from parent node[poids,below] {$0{,}95$}}
    edge from parent node[poids,below] {$0{,}6$}};
\end{tikzpicture}
\end{center}
\begin{calculs}
P(B\cap R) &= P(B)\times P_B(R) &\qquad& \rem{chemin $B$ puis $R$}\\
P(B\cap R) &= 0{,}4\times 0{,}15 &&\\
P(B\cap R) &= \res{0{,}06} &&
\end{calculs}
\explique{$R$ est réalisé par deux chemins : $B$ puis $R$, et $\overline{B}$ puis
$R$.}{on cherche tous les chemins}
\begin{calculs}
P(R) &= P(B\cap R)+P\big(\overline{B}\cap R\big) &\qquad& \rem{somme des deux chemins}\\
P(R) &= 0{,}06+0{,}6\times 0{,}05 &&\\
P(R) &= 0{,}06+0{,}03 &&\\
P(R) &= \res{0{,}09} &&
\end{calculs}
\explique{$9\,\%$ des élèves arrivent en retard au moins une fois par mois.}{on
conclut par une phrase}
\cas{b)}
\explique{Au premier tirage : $4$ rouges sur $6$ jetons. Après un rouge, il reste $3$
rouges et $2$ jaunes ; après un jaune, $4$ rouges et $1$ jaune.}{on garde le jeton :
il en reste $5$}
\begin{center}
\begin{tikzpicture}[arbre]
\node {}
  child {node {$R$}
    child {node {$R$} edge from parent node[poids,above] {$\frac{3}{5}$}}
    child {node {$J$} edge from parent node[poids,below] {$\frac{2}{5}$}}
    edge from parent node[poids,above] {$\frac{4}{6}$}}
  child {node {$J$}
    child {node {$R$} edge from parent node[poids,above] {$\frac{4}{5}$}}
    child {node {$J$} edge from parent node[poids,below] {$\frac{1}{5}$}}
    edge from parent node[poids,below] {$\frac{2}{6}$}};
\end{tikzpicture}
\end{center}
\begin{calculs}
P(\text{deux rouges}) &= \dfrac{4}{6}\times\dfrac{3}{5} &\qquad& \rem{chemin $R$ puis $R$}\\
P(\text{deux rouges}) &= \dfrac{12}{30} &&\\
P(\text{deux rouges}) &= \dfrac{2\times\barre{6}}{5\times\barre{6}} && \rem{on simplifie par $6$}\\
P(\text{deux rouges}) &= \res{\dfrac{2}{5}} &&
\end{calculs}
\explique{« Même couleur » : les chemins $R$ puis $R$, et $J$ puis $J$.}{deux chemins}
\begin{calculs}
P(\text{même couleur}) &= \dfrac{12}{30}+\dfrac{2}{6}\times\dfrac{1}{5} &\qquad& \rem{somme des deux chemins}\\
P(\text{même couleur}) &= \dfrac{12}{30}+\dfrac{2}{30} &&\\
P(\text{même couleur}) &= \dfrac{14}{30} &&\\
P(\text{même couleur}) &= \res{\dfrac{7}{15}} && \rem{on simplifie par $2$}
\end{calculs}
\end{correction}

\etape{À vous de jouer}
\begin{questions}[label=\alph*)]
  \item Dans une salle de sport, $30\,\%$ des abonnés viennent le matin, les autres
    le soir. $20\,\%$ des abonnés du matin suivent un cours collectif, contre
    $50\,\%$ des abonnés du soir. On choisit un abonné au hasard ; on note $M$ :
    « l'abonné vient le matin » et $C$ : « l'abonné suit un cours collectif ».
    Traduire les données par des probabilités, construire l'arbre pondéré, puis
    calculer $P(M\cap C)$ et $P(C)$.
  \item Un sac contient $5$ jetons verts et $2$ jetons bleus, indiscernables au
    toucher. On tire un jeton au hasard, on le garde, puis on tire un second jeton.
    Construire l'arbre pondéré, puis calculer la probabilité d'obtenir deux jetons
    verts, et celle d'obtenir deux jetons de la même couleur.
\end{questions}
\reponse{a) $P(M\cap C)=0{,}06$ ; $P(C)=0{,}41$ ; b)~$\frac{10}{21}$ ; $\frac{11}{21}$.}
\end{methodebox}

\afaire{exercices 23 à 26}

\partiecours{Loi des grands nombres et simulation}

\begin{definitionbox}[Définition 10 --- Échantillon, fréquence observée]
On répète $n$ fois la même expérience aléatoire, de façon indépendante : les $n$
résultats forment un \textbf{échantillon de taille $n$}. La \textbf{fréquence
observée} d'un événement $A$ est
\[
f=\frac{\text{nombre de fois où }A\text{ est réalisé}}{n}
\]
\end{definitionbox}

\begin{proprietebox}[Propriété 8 --- Loi des grands nombres]
Lorsque $n$ est grand, sauf exception, la fréquence observée est proche de la
probabilité.

\smallskip
On peut donc estimer une probabilité inconnue (dé truqué, punaise) par la fréquence
observée sur un grand échantillon.

\smallskip
\textbf{Ordre de grandeur de l'erreur} (admis) : dans environ $95\,\%$ des
échantillons de taille $n$, l'écart entre la fréquence $f$ et la probabilité $p$ est
inférieur à $\dfrac{1}{\sqrt{n}}$.
\begin{center}
\renewcommand{\arraystretch}{2.4}
\begin{tabular}{|l|c|c|c|}
\hline
Taille $n$ & $100$ & $10\,000$ & $1\,000\,000$ \\
\hline
Écart inférieur à $\dfrac{1}{\sqrt{n}}$ & $0{,}1$ & $0{,}01$ & $0{,}001$ \\
\hline
\end{tabular}
\end{center}
Pour diviser l'erreur par $10$, il faut $100$ fois plus d'essais.
\end{proprietebox}

\begin{remarquebox}[En Python]
Après \texttt{from random import randint, random} :
\begin{itemize}
  \item \texttt{randint(1, 6)} donne un nombre entier au hasard entre $1$ et $6$
    (inclus) : c'est un lancer de dé équilibré ;
  \item \texttt{random()} donne un nombre au hasard entre $0$ et $1$ ; le test
    \texttt{random() < 0.7} est vrai avec la probabilité $0{,}7$.
\end{itemize}
\end{remarquebox}

\begin{methodebox}[Méthode 9 --- Simuler une expérience en Python]
\etapeexemples
\begin{questions}[label=\alph*)]
  \item On lance un dé équilibré à six faces ; on gagne si on obtient $5$ ou $6$.
    Calculer la probabilité $p$ de gagner. Expliquer ce que fait la fonction
    ci-dessous, puis ce que montrent les résultats obtenus. Pour $n=100$, puis pour
    $n=1\,000$, l'écart entre la fréquence et $p$ est-il inférieur à
    $\dfrac{1}{\sqrt{n}}$ ?
\end{questions}
\begin{python}
from random import randint

def frequence_gain(n):
    gains = 0
    for k in range(n):
        de = randint(1, 6)
        if de >= 5:
            gains = gains + 1
    return gains / n
\end{python}
\begin{center}
\renewcommand{\arraystretch}{1.5}
\begin{tabular}{|l|c|c|c|c|}
\hline
$n$ & $10$ & $100$ & $1\,000$ & $100\,000$ \\
\hline
\texttt{frequence\_gain(n)} & $0{,}5$ & $0{,}29$ & $0{,}343$ & $0{,}33281$ \\
\hline
\end{tabular}
\end{center}
\begin{questions}[label=\alph*),start=2]
  \item Une joueuse de basket réussit un lancer franc avec la probabilité $0{,}7$.
    Compléter les deux fonctions ci-dessous : \texttt{lancer()} renvoie $1$ si le
    lancer est réussi, $0$ sinon ; \texttt{frequence(n)} renvoie la fréquence des
    lancers réussis sur $n$ lancers. De quelle valeur \texttt{frequence(100000)}
    sera-t-il proche ?
\end{questions}
\begin{python}
from random import random

def lancer():
    if random() < .....:
        return 1
    else:
        return .....

def frequence(n):
    reussis = 0
    for k in range(n):
        reussis = reussis + lancer()
    return .....
\end{python}
\begin{correction}
\cas{a)}
\explique{On gagne avec $5$ ou $6$, donc $p=\dfrac{2}{6}=\res{\dfrac{1}{3}}$, soit environ
$0{,}333$.}{équiprobabilité sur les $6$ faces}
\explique{La boucle simule $n$ lancers ; \texttt{gains} compte les lancers gagnants,
et la fonction renvoie \res{\texttt{gains / n}, la fréquence observée des gains}.}{nombre
de gains divisé par $n$}
\explique{Plus $n$ est grand, plus la fréquence est proche de $\dfrac{1}{3}$ ; avec
$n=10$, elle peut en être loin. C'est \res{la loi des grands nombres}.}{propriété 8}
\explique{Pour $n=100$ : l'écart vaut environ $0{,}333-0{,}29$, soit $0{,}043$, et
$\dfrac{1}{\sqrt{100}}=0{,}1$ : \res{l'écart est inférieur à $0{,}1$}.}{$\sqrt{100}=10$}
\explique{Pour $n=1\,000$ : l'écart vaut environ $0{,}343-0{,}333$, soit $0{,}01$, et
$\dfrac{1}{\sqrt{1\,000}}\approx 0{,}032$ : \res{l'écart est inférieur à
$0{,}032$}.}{calculatrice pour $\sqrt{1\,000}$}
\cas{b)}
\begin{python}
from random import random

def lancer():
    if random() < 0.7:
        return 1
    else:
        return 0

def frequence(n):
    reussis = 0
    for k in range(n):
        reussis = reussis + lancer()
    return reussis / n
\end{python}
\explique{\texttt{random() < 0.7} est vrai avec la probabilité $0{,}7$ ; on renvoie $0$
pour un lancer manqué, et \texttt{reussis / n} est la fréquence.}{on additionne des $1$
et des $0$}
\explique{D'après la loi des grands nombres, \texttt{frequence(100000)} sera proche de
\res{$0{,}7$}.}{$n=100\,000$ est grand}
\end{correction}

\etape{À vous de jouer}
\begin{questions}[label=\alph*)]
  \item On lance un dé équilibré à six faces ; on gagne si on obtient $6$. Calculer la
    probabilité $p$ de gagner. Expliquer ce que fait la fonction ci-dessous, puis ce
    que montrent les résultats obtenus. Pour $n=100$, puis pour $n=1\,000$, l'écart
    entre la fréquence et $p$ est-il inférieur à $\dfrac{1}{\sqrt{n}}$ ?
\end{questions}
\begin{python}
from random import randint

def frequence_six(n):
    six = 0
    for k in range(n):
        de = randint(1, 6)
        if de == 6:
            six = six + 1
    return six / n
\end{python}
\begin{center}
\renewcommand{\arraystretch}{1.5}
\begin{tabular}{|l|c|c|c|c|}
\hline
$n$ & $10$ & $100$ & $1\,000$ & $100\,000$ \\
\hline
\texttt{frequence\_six(n)} & $0{,}3$ & $0{,}19$ & $0{,}172$ & $0{,}16641$ \\
\hline
\end{tabular}
\end{center}
\begin{questions}[label=\alph*),start=2]
  \item Un archer atteint la cible avec la probabilité $0{,}4$. Compléter les deux
    fonctions ci-dessous : \texttt{tir()} renvoie $1$ si la cible est atteinte, $0$
    sinon ; \texttt{frequence(n)} renvoie la fréquence des tirs réussis sur $n$ tirs.
    De quelle valeur \texttt{frequence(100000)} sera-t-il proche ?
\end{questions}
\begin{python}
from random import random

def tir():
    if random() < .....:
        return 1
    else:
        return .....

def frequence(n):
    reussis = 0
    for k in range(n):
        reussis = reussis + tir()
    return .....
\end{python}
\reponse{a) $p=\frac{1}{6}\approx 0{,}167$ ; les fréquences se rapprochent de
$\frac{1}{6}$ ; écarts d'environ $0{,}023$ ($<0{,}1$) et $0{,}005$ ($<0{,}032$) ; b) \texttt{random() < 0.4}, \texttt{return 0},
\texttt{return reussis / n} ; proche de $0{,}4$.}
\end{methodebox}

\afaire{exercices 27 à 29}

\end{document}
